\section{Number-theory consequences with the hypotheses separated}
\label{sec:number-theory}

Throughout this section $A$ denotes the user-designated Alp\"oge theorem
for the Riemann zeta function and Dirichlet L-functions.  It is an adopted
cited input, not a claim that this run rebuilt the unseen Lean dependencies.
$M$ denotes the conditional all-number-field unitary Hecke-GRH assertion
of Masullo, \S1.2.  Every use of $M$ below remains conditional.  The
accompanying catalogue is \texttt{work/number\_theory/rows.json}; its
literature records include retrieval evidence and distinguish a checked
statement from a reread of an entire classical proof.

There are three different geometric translations.  Rational primes are
the closed points of $\operatorname{Spec}\mathbb Z$; prime ideals are the
closed points of $\operatorname{Spec}\mathcal O_K$; and a conjugacy condition
on Frobenius is a condition on a finite \'{e}tale cover of an open arithmetic
curve.  The last dictionary is the Galois-category equivalence in SGA~1,
Expos\'e~V, \S7.  The class group is
$\operatorname{Pic}(\operatorname{Spec}\mathcal O_K)$ and its class fields
are finite \'{e}tale covers.  These translations identify the objects to
which the estimates apply.  They do not turn number-field GRH into a
consequence of finite-field purity.  The inherited cote-1 trace programme
is a comparison of explicit formulas; no new archive scan is certified by
this section.  Exact archive work belongs to the geometry section.

\subsection{The classical arithmetic bounds being used}

\paragraph{NT-P01--NT-P04: rational primes and gaps.}
Schoenfeld, \emph{Sharper bounds for the Chebyshev functions
$\theta(x)$ and $\psi(x)$. II} (1976), Theorem~10, printed p.~337,
equations (6.2) and (6.3), gives, under RH,
\[
 |\psi(x)-x|<\frac{\sqrt{x}\log^2x}{8\pi}\quad(x\ge73.2),\qquad
 |\theta(x)-x|<\frac{\sqrt{x}\log^2x}{8\pi}\quad(x\ge599).
\]
Equation (6.1) replaces $\log^2x$ by $(\log x-2)\log x$ for
$x\ge23\cdot10^8$.  Corollary~1 gives
$|\pi(x)-\operatorname{li}(x)|<\sqrt{x}\log x/(8\pi)$ for
$x\ge2657$.  Here $\operatorname{li}$ has its usual normalization.
The constants in Theorem~10 were checked in the primary indexed text.
An independently retrieved primary proof locator for the same
prime-counting bound is Johnston (2022), Theorem~3.2, equation (3.6),
\emph{The Ramanujan Journal} \textbf{59}, 1307--1321.
It requires RH only to a height $T$ satisfying
$9.06\sqrt{x/\log x}/\log\log x\le T$; under $A$ one may choose
such a $T$ for each $x$.  Its \S2 explicitly says its main theorem applies
also to the ordinary, non-half-weight counting functions.  Equation (3.3)
there also improves the threshold for the displayed $\psi$ inequality
to $x\ge59$.  Thus the inaccessible original Corollary~1 page is no
longer a proof-source gate for NT-P03.

Dudek--Greni\'e--Molteni, \emph{Primes in explicit short intervals on RH}
(2016), Theorem~1.1, author PDF p.~1, gives a prime in
\[
 (x-c\sqrt{x}\log x,\ x+c\sqrt{x}\log x),
 \qquad c=\tfrac12+2/\log x,\qquad x\ge2.
\]
The same theorem gives at least $\sqrt{x}$ primes after replacing $c$ by
$c+1$.  These results imply the classical
$p_{n+1}-p_n=O(\sqrt{p_n}\log p_n)$ consequence of RH.  Neither this
bound nor the assertion of RH supplies a proof here of Cram\'er's
$O(\log^2p_n)$ conjecture, twin primes, or the binary Goldbach conjecture.
Those problems require additional information about prime correlations.

\paragraph{NT-P05--NT-P08: progressions and nonresidues.}
Greni\'e--Molteni--Perelli (2016), Theorem~1, author PDF p.~2, proves that
absolute $x_0,c_1,c_2>0$ work in
\[
 \pi(x+h;q,a)-\pi(x;q,a)\ge
 c_2\frac{h}{\varphi(q)\log x},\qquad
 c_1\varphi(q)\sqrt{x}\log x\le h\le x,\quad x\ge x_0,
\]
assuming the Dirichlet GRH modulo $q$ and $(a,q)=1$.  The following
sentence in that source gives the asymptotic if the ratio to the lower
scale tends to infinity.  We have not converted the unspecified absolute
constants into numerical constants.

Lamzouri--Li--Soundararajan (LLS), corrected
\texttt{arXiv:1309.3595v3}, is used for the sharp explicit conductor
statements.  Corollary~1.2, author p.~3, gives the least prime in a reduced
class modulo $q>3$ at most $(\varphi(q)\log q)^2$.
Corollary~1.1, p.~2, gives a quadratic nonresidue below $(\log q)^2$ for
prime $q\ge5$.  Theorem~1.1(1), p.~2, gives the proper-subgroup bound
recorded with its explicit functions $A(q),B(q)$ in NT-P08.
These are the already carried IA02 consequences, with the corrected
version and hypotheses rechecked; they are not new findings of this run.
Geometrically the progression is a Frobenius class in a cyclotomic cover,
and the nonresidue is detected by a finite rank-one character.

\paragraph{NT-P09: the least primitive root is a different quantifier.}
Shoup, \emph{Searching for primitive roots in finite fields} (1992),
Theorem~3, author PDF p.~4, gives under Dirichlet GRH
\[
 g(p)\ll r^4(\log r+1)^4(\log p)^2,
 \qquad r=\omega(p-1).
\]
The printed theorem was checked on the page image.  Here is the short
simplification to $g(p)\ll\log^6p$.  If $q_1<\cdots<q_r$ are the distinct
prime divisors of $p-1$, then $q_j\ge j+1$ and therefore
$p-1\ge\prod_{j=1}^r(j+1)$.  For $r\ge4$ the last at least $r/2$ factors
are at least $r/2$, so
$\log(p-1)\ge(r/2)\log(r/2)$.  Consequently
$r(\log r+1)\ll\log p$; enlarging the absolute constant handles $r<4$.
Substitution proves the assertion.  This finds a small generator for a
fixed $\mathbb F_p^\times$.  The density theorem for one fixed integer
as $p$ varies uses different, generally nonabelian, fields (NT-A01).

\subsection{A quantitative Frobenius trace transfer}

\paragraph{NT-C01 and NT-C02: the precise cited Chebotarev input.}
Let $L/K$ be a finite Galois extension, $G=\operatorname{Gal}(L/K)$,
$n_L=[L:\mathbb Q]$, and $\Delta_L=|\operatorname{disc}L|$.
For a conjugacy class $C$ put $\delta_C=|C|/|G|$.  For unramified
primes only, define
\[
 \psi_C(x)=\sum_{\substack{\mathfrak p\text{ unramified},\ m\ge1\\
                    (N\mathfrak p)^m\le x}}
 \mathbf1_C(\operatorname{Frob}_{\mathfrak p}^{m})\log N\mathfrak p,
 \qquad
 \pi_C(x)=\sum_{\substack{\mathfrak p\text{ unramified}\\N\mathfrak p\le x}}
             \mathbf1_C(\operatorname{Frob}_{\mathfrak p}).
\]
Set
\begin{align*}
 E_L(x)&=\sqrt{x}\left[
       \left(\frac{\log x}{2\pi}+2\right)\log\Delta_L+
       \left(\frac{\log^2x}{8\pi}+2\right)n_L\right],\qquad x\ge1,\\
 B_L(x)&=\sqrt{x}\left[
       \left(\frac1{2\pi}+\frac3{\log x}\right)\log\Delta_L+
       \left(\frac{\log x}{8\pi}+\frac1{4\pi}+\frac6{\log x}\right)n_L
                         \right],\qquad x\ge2,\\
 \operatorname{Li}_2(x)&=\int_2^x\frac{dt}{\log t}.
\end{align*}
Greni\'e--Molteni, \emph{An explicit Chebotarev density theorem under GRH}
(2019), Theorem~1.1 and Corollary~1.2, author PDF pp.~2--3, give
\[
 |\psi_C(x)-\delta_Cx|\le\delta_C E_L(x),\qquad
 |\pi_C(x)-\delta_C\operatorname{Li}_2(x)|\le\delta_C B_L(x),
\tag{NT-C}
\]
under RH for $\zeta_L$.  The theorem explicitly includes the unramified
version of $\psi_C$.  Its constants and ranges were checked on the page
images.  In particular Artin holomorphy is not an additional hypothesis.
$M$ supplies this input for every $L$; $A$ supplies it if $L/\mathbb Q$
is abelian by the factorization into primitive Dirichlet functions.

\paragraph{NT-F01. Statement and full proof.}
For a complex class function $f$ on $G$ put
\[
 \mu_f=\frac1{|G|}\sum_{g\in G}f(g),\qquad
 \|f\|_{1,G}=\frac1{|G|}\sum_{g\in G}|f(g)|,
\]
and define $\Psi_f$ and $\Pi_f$ by replacing $\mathbf1_C$ in the two
sums above by $f$.  Then
\[
 |\Psi_f(x)-\mu_fx|\le\|f\|_{1,G}E_L(x),\qquad
 |\Pi_f(x)-\mu_f\operatorname{Li}_2(x)|\le\|f\|_{1,G}B_L(x).
\tag{NT-F01}
\]
If $f=\chi_V$ is the character of a finite-dimensional complex
representation, then $\mu_f=\dim V^G$ and
$\|f\|_{1,G}\le\sqrt{\langle\chi_V,\chi_V\rangle_G}$.

To prove the estimates, choose one representative $g_C$ from each
conjugacy class.  The finite identities
\[
 f=\sum_C f(g_C)\mathbf1_C,
 \qquad \mu_f=\sum_C\delta_Cf(g_C)
\]
give $\Psi_f-\mu_fx=\sum_Cf(g_C)(\psi_C-\delta_Cx)$.
The triangle inequality and (NT-C) bound its modulus by
$E_L(x)\sum_C\delta_C|f(g_C)|=E_L(x)\|f\|_{1,G}$.
The same finite argument proves the assertion for $\Pi_f$.
The averaging operator $|G|^{-1}\sum_g\rho(g)$ is the projection to $V^G$,
so its trace is $\mu_{\chi_V}=\dim V^G$.  Cauchy--Schwarz for the
probability measure assigning mass $1/|G|$ to each group element yields
\[
 \left(\frac1{|G|}\sum_g|\chi_V(g)|\right)^2
 \le\frac1{|G|}\sum_g|\chi_V(g)|^2
 =\langle\chi_V,\chi_V\rangle_G.
\]
This proves every assertion.  This is a linear functional estimate on the
complexified representation ring; virtual differences need not be
nonnegative, so their error is controlled by the displayed norm.

\paragraph{NT-F02. Finite \texorpdfstring{\'{e}}{e}tale point counts, with the boundary.}
Let $U$ be a nonempty open subscheme of
$\operatorname{Spec}\mathcal O_K$, and let $X\to U$ be finite \'{e}tale
of degree $d$.  Let $S_X$ be the set of geometric points of its generic
fibre, let $L/K$ be its finite Galois splitting field, and let $r$ be the
number of $G$-orbits on $S_X$, equivalently the number of generic connected
components.  Let $S$ be the finite set of primes not in $U$ which are
unramified in $L$.  Under RH for $\zeta_L$,
\[
 \left|\sum_{\substack{\mathfrak p\in U,\ m\ge1\\(N\mathfrak p)^m\le x}}
       \#X(k_{\mathfrak p^m})\log N\mathfrak p-rx\right|
 \le rE_L(x)+d|S|\log x,\qquad x\ge1.
\tag{NT-F02}
\]
Here $k_{\mathfrak p^m}$ denotes the degree-$m$ extension of
$k_{\mathfrak p}$.  In particular, when no additional unramified prime is
removed, the error is $rE_L(x)$, independent of the degree $d$.

For the proof, finite \'{e}taleness implies that inertia at a prime of $U$
acts trivially on the generic geometric fibre.  The action of
$\operatorname{Gal}(L/K)$ on this fibre is faithful by the definition of
the splitting field, so $L/K$ is unramified at every prime of $U$.
Let $\chi(g)=\#S_X^g$ be the permutation character.  The elementary
Frobenius fixed-point formula for a finite \'{e}tale scheme over a finite
field is
$\#X(k_{\mathfrak p^m})=\chi(\operatorname{Frob}_{\mathfrak p}^{m})$:
its geometric points rational over that field are precisely the points
fixed by the $m$th Frobenius power.  Also $0\le\chi(g)\le d$.
Double counting pairs $(g,s)$ with $gs=s$ gives
\[
 \sum_{g\in G}\chi(g)=\sum_{s\in S_X}|G_s|=r|G|.
\]
Thus both $\mu_\chi$ and $\|\chi\|_{1,G}$ equal $r$.  NT-F01 bounds
the sum over every unramified prime by $rx$ with error $rE_L(x)$.
For each omitted prime, the nonnegative lost contribution is at most
\[
 d\left\lfloor\frac{\log x}{\log N\mathfrak p}\right\rfloor
          \log N\mathfrak p\le d\log x.
\]
If $N\mathfrak p>x$ the contribution is zero.  Summing this bound over
$S$ and applying the triangle inequality proves (NT-F02), including
$x=1$.  Notice that replacing the nonnegative character norm by the square
root of its second moment would give a weaker bound; the orbit count is
the relevant parameter here.

\paragraph{NT-F03. Splitting types.}
For a partition $\lambda$ of $d=|S_X|$, let $D_\lambda\subseteq G$ be
the elements with cycle partition $\lambda$ on $S_X$, and put
$\delta_\lambda=|D_\lambda|/|G|$.  This is a union of conjugacy classes.
Sum the second inequality (NT-C) over these classes and use the triangle
inequality.  Their densities sum to $\delta_\lambda$, proving
\[
 \left|\#\{\mathfrak p:\ N\mathfrak p\le x,\ \mathfrak p
 \text{ unramified},\ \operatorname{Frob}_{\mathfrak p}\in D_\lambda\}
 -\delta_\lambda\operatorname{Li}_2(x)\right|
 \le\delta_\lambda B_L(x).
\]
If $D_\lambda$ is empty both sides before the absolute value are zero.
For the canonical finite \'{e}tale model of the generic algebra over the
unramified locus, each Frobenius orbit is one closed point of its fibre,
and its orbit length is that point's residue degree.  Hence this cycle
partition is exactly its splitting type.  For the originally specified
model $X$ over $U$, restrict the counted primes to $U$; at most $|S|$
terms have been removed, so the same estimate holds with the additional
error $|S|$.  This proves both interpretations without assigning a fibre
of $X$ outside its specified base.

\paragraph{NT-F04. Avoiding a prescribed finite set.}
Let $D\subseteq G$ be a nonempty conjugacy-stable set of density
$\delta_D$, and let $S$ be any finite set of primes of $K$.  The same
finite-union argument gives
$\pi_D(x)\ge\delta_D(\operatorname{Li}_2(x)-B_L(x))$.
At most $|S|$ of the counted primes belong to $S$.  If
\[
 \delta_D\bigl(\operatorname{Li}_2(x)-B_L(x)\bigr)>|S|,
\]
the count outside $S$ is strictly positive.  It is an integer, so an
eligible prime exists.  This proves the avoidance assertion without adding
an unsupported arbitrary-exclusion clause to Bach--Sorenson's least-prime
theorem.  That theorem is NT-C03, Theorem~5.1, printed p.~1729, with
RH hypotheses specified in its \S2, p.~1720.

\subsection{Class numbers, Kummer towers, and Euclidean algorithms}

\paragraph{NT-H01--NT-H04: the regulator normalization.}
Write
\begin{align*}
 U(q)&=2e^\gamma\left(\log\log q-\log2+\tfrac12+
                                      \frac1{\log\log q}\right),\\
 V(q)&=\frac{12e^\gamma}{\pi^2}\left(\log\log q-\log2+\tfrac12+
                   \frac1{\log\log q}+\frac{14\log\log q}{\log q}\right).
\end{align*}
LLS Theorem~1.5, author pp.~3--4, gives
$V(q)^{-1}\le|L(1,\chi)|\le U(q)$ for primitive conductor
$q\ge10^{10}$.  Its Corollary~1.3 is the already carried imaginary
quadratic bound in NT-H02.  The inequality is not a two-sided estimate of
order $\log\log q$: the reciprocal lower bound must be retained.

Here is the full deduction of NT-H03, including its normalization.
For $D>0$ fundamental put $K=\mathbb Q(\sqrt D)$ and let $R_D$ be the
regulator of the full unit group modulo its torsion, so
$R_D=\log\varepsilon$ for a fundamental unit $\varepsilon>1$, regardless
of its norm's sign.  The analytic class-number formula is
\[
 \mathop{\rm Res}_{s=1}\zeta_K(s)=\frac{2h_DR_D}{\sqrt D}.
\]
An exact independent locator is Milne, \emph{Class Field Theory}, v4.03,
Chapter~VI, Corollary~2.12 and Example~2.13, printed p.~191.
Since $\zeta_K(s)=\zeta(s)L(s,\chi_D)$ and $\zeta$ has residue one,
$L(1,\chi_D)=2h_DR_D/\sqrt D>0$.
Multiplying the two inequalities in NT-H01 by $\sqrt D/2$ proves
\[
 \frac{\sqrt D}{2V(D)}\le h_DR_D\le\frac{\sqrt D}{2}U(D),
 \qquad D\ge10^{10}.
\]
This completes the short proof; no bound on $R_D$ was suppressed.

For completeness, the already inherited abelian-field formulation
NT-H04 is the following precise product calculation.  If $K/\mathbb Q$
is abelian with Dirichlet-character group $X_K$, then
\[
 h_KR_K=\frac{w_K\sqrt{\Delta_K}}{2^{r_1}(2\pi)^{r_2}}
              \prod_{\substack{\chi\in X_K\\\chi\ne1}}L(1,\chi).
\tag{NT-H04}
\]
Indeed Artin factorization gives the product for $\zeta_K/\zeta$,
and taking the residue at one gives this formula by the same class-number
formula.  Its left side and the prefactor are positive, so taking moduli
turns the product into $\prod|L(1,\chi)|$.  For any factors whose
conductors are at least $10^{10}$ insert $U$ and $1/V$; retain the finite
remaining product exactly or bound it separately.  This proves every
claimed product bound.  Dividing by an unestimated regulator does not
prove the corresponding bound on $h_K$.  This is the correction already
made in RETURN\_ZETA\_TWO IA02 and IA09, not a new result of this run.

\paragraph{NT-A01 and NT-A02: the nonabelian input in Artin's primitive-root problem.}
For a fixed integer $g\ne-1$ that is not a square, put
$K_n=\mathbb Q(\zeta_n,g^{1/n})$.  Hooley's classical theorem gives a
positive density
\[
 \delta(g)=\sum_{n\ge1}\frac{\mu(n)}{[K_n:\mathbb Q]}
\]
of primes where $g$ is a primitive root, assuming RH for these Kummer
fields.  We cite Hooley, \emph{On Artin's conjecture} (1967), the
unnumbered theorem on pp.~219--220.  The original internal page has not
been freshly obtained here; a fully inspected precise proof locator for
the density implication is Lenstra, \emph{On Artin's conjecture and
Euclid's algorithm in global fields} (1977), Theorem~(3.1), printed
p.~206, specialized to $K=F=\mathbb Q$, $W=\langle g\rangle$, $k=1$,
with density (2.15), p.~206, and positivity criterion (4.1), p.~207.
The series uses actual field degrees; replacing it by the universal
Artin constant for every $g$ loses the entanglement correction.

Lenstra's version also imposes a conjugacy-stable Frobenius condition
$C\subseteq\operatorname{Gal}(F/K)$ and asks that the residue image of
a finitely generated positive-rank $W\subseteq K^\times$ have index
dividing $k$.  Its auxiliary fields are formed from
$K(\zeta_{q(\ell)},W^{1/q(\ell)})$ with
$q(\ell)=\ell^{v_\ell(k)+1}$, taking the compositum for primes dividing
a squarefree $n$.  Theorem~(3.1) assumes RH for these fields' zeta
functions.  Theorem~(4.1) gives the finite obstruction criterion for
positive density.  $M$ supplies the hypotheses; $A$ alone does not
supply them because the Kummer fields need not be abelian over
$\mathbb Q$.  In geometric terms this is a tower of Kummer torsors of
$\mathbb G_m$, not merely the cyclotomic torsion cover.

\paragraph{NT-A03: a PID becomes Euclidean under the stated Kummer GRH.}
Lenstra's Theorem~(9.1), printed p.~218, states the result recorded in
NT-A03 for the ring of $S$-integers with $|S|\ge2$ and trivial ideal
class group.  Equation (9.2) gives its least Euclidean function.
The original page image confirms the weak inequality $|S|\ge2$;
the OCR's apparent strict inequality was not adopted.  The result is a
classical cited theorem with an exact proof locator, not a new proof of
the Euclidean algorithm.  The independent PID hypothesis and the
distinction from a norm-Euclidean algorithm remain in the catalogue.

\subsection{Growth, cancellation, and zero statistics}

\paragraph{NT-L02 and NT-L03. The fixed-character Lindel\"of deduction.}
The classical input is Simoni\v c, \emph{Estimates for L-functions in the
critical strip under GRH with effective applications} (2023), Theorem~1,
equations (4)--(6), proved in its \S3.  For a function $F$ in the Selberg
class with polynomial Euler product, zero-free to the right of the
critical line, it implies, for fixed $\delta>0$ and large $|t|$,
\[
 \log|F(\sigma+it)|,\qquad -\log|F(\sigma+it)|
 \ll_{F,\delta}(\log|t|)^{1-2\delta}\log\log|t|,
 \quad \tfrac12+\delta\le\sigma\le1.
\tag{NT-L-log}
\]
For $\delta\ge1/2$ one uses instead a smaller positive $\delta$ and
the absolutely convergent Euler product on $\sigma>1$.
The modulus bound for $\log F$ in the cited theorem gives both signs of
its real part.  Dividing the right side of (NT-L-log) by $\log|t|$
shows that it tends to zero for fixed $0<\delta<1/2$.  Hence, for every
$\epsilon>0$,
\[
 F(\sigma+it),\quad F(\sigma+it)^{-1}
 \ll_{F,\delta,\epsilon}(1+|t|)^\epsilon
 \qquad(\sigma\ge\tfrac12+\delta),
\tag{NT-L-right}
\]
with pole neighbourhoods omitted if $F$ has a pole.  The finite range of
$t$ is absorbed into the constant on any closed set avoiding poles and
zeros.  On $\sigma\ge1+\delta$ the Euler product gives the required
uniform bound directly; compact intermediate strips are already covered
by (NT-L-log) or the cited theorem's adjacent range.

Here is the passage to the critical line, often suppressed in invoking
Lindel\"of.  Fix $0<\delta<1/2$.  If $F$ has degree $d$, its functional
equation and Stirling's formula imply
\[
 |F(\tfrac12-\delta+it)|
 \ll_{F,\delta}(1+|t|)^{d\delta}
       |F^\vee(\tfrac12+\delta-it)|.
\]
The dual has the same zero-free hypothesis, so NT-L-right bounds this by
$O((1+|t|)^{d\delta+\eta})$ for every $\eta>0$.  The right boundary
$\sigma=1/2+\delta$ has exponent $\eta$.
The Phragm\'en--Lindel\"of principle on this strip gives at its midpoint
\[
 |F(\tfrac12+it)|\ll_{F,\delta,\eta}
                       (1+|t|)^{d\delta/2+\eta}.
\]
Its hypotheses hold because the standard functional equation and
finite-order analytic continuation give polynomial growth in each fixed
vertical strip; if there is a pole, apply the principle above and below
a fixed height containing it.  Choose $\delta$ and $\eta$ so that
$d\delta/2+\eta<\epsilon$.  The remaining bounded interval is compact,
and the critical line contains no pole for these L-functions.  This
proves $F(1/2+it)\ll_{F,\epsilon}(1+|t|)^\epsilon$.

For primitive Dirichlet functions the hypotheses are classical, and $A$
is the zero-free input; this proves NT-L02.  For a fixed non-pure-norm
unitary Hecke character, the Hecke L-function is entire of finite order,
its gamma shifts have nonnegative real parts, and its degree is
$[K:\mathbb Q]$.  On grouping its ideal Euler factors by rational primes,
one obtains a polynomial Euler product of that degree with local roots
of modulus at most one.  The Dirichlet coefficients are bounded by
$d_{[K:\mathbb Q]}(n)=O_\epsilon(n^\epsilon)$: each local coefficient
is bounded by the corresponding coefficient of $(1-p^{-s})^{-[K:\mathbb Q]}$.
Thus it meets the explicit axioms in Simoni\v c \S2.1, and $M$ supplies
zero-freeness.  If the character is a pure unitary norm power, its
L-function is a vertical translate of $\zeta_K$, to which the same
argument applies before translation.  Imprimitive Euler factors, when
included in the convention, form a fixed finite product; they and their
inverses are bounded on each $\sigma\ge1/2+\delta$, since their roots
have modulus at most one.  This proves NT-L03 in all the conventions just
specified.  Constants may depend on the field, character, infinity type,
and $\epsilon$; this proof makes no uniform conductor-aspect claim.

\paragraph{NT-L05 and NT-L06. Mertens, Liouville, and reciprocal tails.}
Titchmarsh, second edition, Theorem~14.25(C), printed p.~370, is the exact
classical equivalence
$\mathrm{RH}\Longleftrightarrow M(x)=O_\epsilon(x^{1/2+\epsilon})$
for every $\epsilon>0$.  We use that theorem rather than reproduce its
contour proof.  The following short deductions are included in full.
For a prime power $p^a$,
\[
 \sum_{d^2\mid p^a}\mu(p^a/d^2)=
 \begin{cases}1,&a\text{ even},\\-1,&a\text{ odd},\end{cases}
\]
because only the term with remaining exponent zero or one is nonzero.
Both sides are multiplicative, hence
$\lambda(n)=\sum_{d^2\mid n}\mu(n/d^2)$ for every $n$.
Consequently
\[
 \left|\sum_{n\le x}\lambda(n)\right|
 \le\sum_{d\le\sqrt x}|M(x/d^2)|
 \ll_\epsilon x^{1/2+\epsilon}
                \sum_{d\ge1}d^{-1-2\epsilon}
 \ll_\epsilon x^{1/2+\epsilon}.
\]
This proves Liouville cancellation, with the dependence on $\epsilon$
retained in the convergent sum.

Let $s=\sigma+it$ and choose $\eta>0$ with $1/2+\eta<\sigma$.
Partial summation for $y>x$ gives
\[
 \sum_{x<n\le y}\frac{\mu(n)}{n^s}
 =M(y)y^{-s}-M(x)x^{-s}
     +s\int_x^y M(u)u^{-s-1}\,du.
\]
The first term tends to zero as $y\to\infty$, and the integral converges
absolutely after using $|M(u)|\ll_\eta u^{1/2+\eta}$.  Its modulus is at
most a constant times
$|s|x^{1/2+\eta-\sigma}/(\sigma-1/2-\eta)$.
Together with the boundary term this proves the stated tail estimate.
It also proves local uniform convergence of the reciprocal Dirichlet
series on $\operatorname{Re}s>1/2$ by choosing $\eta$ below the distance
of a compact set from the boundary.  It agrees with $1/\zeta(s)$ on
$\operatorname{Re}s>1$ by the absolutely convergent Euler product and
therefore throughout that half-plane by analytic continuation.  At
$s=1$, $1/\zeta$ has value zero.  Hence $\sum\mu(n)/n=0$, as asserted.

\paragraph{NT-L04, NT-L07, NT-L08: what moments and zero density say.}
Harper (2013), Theorem~1, author p.~1, gives
$\int_T^{2T}|\zeta(1/2+it)|^{2k}dt\ll_kT(\log T)^{k^2}$ for fixed
$k\ge0$ under RH.  Its proof does not supply the conjectural leading
constant for every $k$.  Simoni\v c (2023), Theorem~2, equations (7),(8),
provides the explicit effective Mertens constants in NT-L07.
For NT-L08, if $\sigma>1/2$, every nontrivial zero covered by $A$ or
conditionally $M$ has real part $1/2<\sigma$, so the indicated counting
set is empty.  At $\sigma=1/2$ this reasoning says nothing about
multiplicity.  Trivial zeros from gamma factors and zeros of removed
finite Euler factors are not silently included in this definition.

\paragraph{NT-Z01 and NT-Z02: the exact pair-correlation range.}
For zeta zero ordinates $0<\gamma,\gamma'\le T$, counted with
multiplicity, put
\[
 F(\alpha,T)=\left(\frac{T}{2\pi}\log T\right)^{-1}
 \sum_{\gamma,\gamma'}T^{i\alpha(\gamma-\gamma')}
                    \frac4{4+(\gamma-\gamma')^2}.
\]
Montgomery (1973), \S1, unnumbered theorem, printed p.~181, proves under
RH the formula in NT-Z01 uniformly on each $0\le\alpha\le1-\epsilon$.
Corollary~2, p.~182, gives the proportion $2/3+o(1)$ of simple zeros.
The extension of the pair-correlation formula beyond that range is
explicitly a conjecture in that paper.  Adopting RH removes the theorem's
hypothesis; it does not replace its range by all $\alpha$.

\paragraph{NT-Z03--NT-Z05: low zeros of the prime-modulus Dirichlet family.}
For prime $q$ average over the $q-2$ nonprincipal characters, with scaled
ordinates $x_{\chi,j}=\gamma_{\chi,j}\log q/(2\pi)$.
Hughes--Rudnick, \emph{Linear statistics of low-lying zeros of
L-functions} (2003), Theorem~3.1, author p.~5, gives the unitary mean for
admissible even test functions with Fourier support in $[-2,2]$.
Theorem~4.1, p.~7, gives variance
$\int_{-1}^1|u|\widehat f(u)^2du$ with support in $[-1,1]$.
Theorem~5.1, p.~9, gives the Gaussian moments only for
$m<2/\alpha$ when the support is in $[-\alpha,\alpha]$.
These identities are unconditional with complex $\gamma$ allowed;
$A$ supplies the real-ordinate interpretation.  They are not an
unrestricted all-moment Gaussian theorem for one fixed test function.

In that source's ordering on p.~1, $\gamma_{\chi,1}$ is the first
\emph{nonnegative} ordinate, so a central zero is included with value zero.
With exactly this convention, Corollary~8.2, p.~20, gives
$\liminf_q\min_{\chi\ne\chi_0}x_{\chi,1}\le1/4$ under $A$.
Theorem~8.3, pp.~20--21, gives a positive proportion for
$x_{\chi,1}<\beta$ if $\beta\ge0.633$.
No strictly-positive-zero variant was inferred by discarding central
zeros.  The geometric comparison is with finite-field family monodromy
and its symmetry group; it does not identify the unrestricted
number-field limiting distribution.

\paragraph{NT-Z06 and NT-Z07: the additional gates.}
No argument in the cited sources proves full GUE, simplicity of every
zero, or central nonvanishing for every quadratic Dirichlet function
from zero location alone.  We record the specific missing correlation
or nonvanishing estimates, without claiming a logical independence
theorem.  For Iwaniec--Luo--Sarnak, \emph{Low lying zeros of families of
L-functions} (2000), the standing assumption on printed p.~56 explicitly
includes RH for the GL$_2$ functions in the family as well as Dirichlet
GRH.  Theorem~1.1, pp.~58--59, gives support $(-2,2)$ for fixed even
weight and squarefree level tending to infinity.  $A$ supplies its
Dirichlet part.  A universal Hecke statement does not identify a general
non-CM GL$_2$ newform with a rank-one Hecke character.  Thus the GL$_2$
hypothesis remains an additional input in this invocation.

\paragraph{NT-Z08 and NT-Z09: two exact classical equivalents.}
Rodgers--Tao, \texttt{arXiv:1801.05914v5}, Theorem~1, author p.~3,
proves $\Lambda\ge0$ for the classical de Bruijn--Newman constant.
The definition and RH equivalence in its \S1 give
$A\Rightarrow\Lambda\le0$, hence $\Lambda=0$.
For the deformation as normalized there, real-rootedness holds exactly
for $t\ge\Lambda$, yielding the two assertions in NT-Z08.
Lagarias, \texttt{arXiv:math/0008177v2}, Theorem~1.1, author p.~1,
is the exact locator for NT-Z09's divisor-sum inequality.
These are cited classical consequences, not additional new proofs of RH.

\subsection{Fixed-character cancellation and the local zero count}

\paragraph{NT-L09 and NT-L10. A full Perron transfer.}
Here is the elementary analytic form used for both rational and ideal
M\"obius sums.  Suppose
$G(s)=\sum_{n\ge1}a_nn^{-s}$ is absolutely convergent for
$\operatorname{Re}s>1$, $|a_n|\le d_D(n)$ for a fixed positive integer
$D$, and $G$ extends holomorphically to $\operatorname{Re}s>1/2$.
Assume that for every $\delta,\eta>0$,
\[
 G(\sigma+it)\ll_{\delta,\eta}(1+|t|)^\eta
 \quad\text{uniformly for }\tfrac12+\delta\le\sigma\le2.
\tag{NT-Perron-input}
\]
Then $\sum_{n\le x}a_n\ll_\epsilon x^{1/2+\epsilon}$ for every
$\epsilon>0$.

We prove the assertion, recording the truncation error as well as the
contour displacement.  It suffices to take $0<\epsilon<1/4$ and
half-integral $x\ge4$, because replacing any real $y$ by
$\lfloor y\rfloor+1/2$ leaves the partial sum unchanged and changes its
size by a bounded factor; the bounded remaining range is absorbed into
the constant.  Put
$c=1+1/\log x$, $T=x^2$, $\delta=\epsilon/4$ and
$\eta=\epsilon/8$.  The classical truncated Perron kernel gives
\[
 \sum_{n\le x}a_n=
 \frac1{2\pi i}\int_{c-iT}^{c+iT}G(s)\frac{x^s}{s}\,ds+R,
\]
with
\[
 |R|\ll\sum_{n\ge1}|a_n|(x/n)^c
           \min\left(1,\frac1{T|\log(x/n)|}\right).
\tag{NT-Perron-error}
\]
This is the usual Perron formula; an exact classical locator is
Titchmarsh, Lemma~3.12 and its proof, printed pp.~61--62.  The kernel
estimate follows by moving the integral of $y^s/s$ to the left or right
according to $y>1$ or $y<1$, taking the residue at zero in the first
case, and integrating the horizontal tails; the alternative trivial
bound handles $|\log y|\le T^{-1}$.

For $n\le x/2$ or $n\ge2x$ the denominator $|\log(x/n)|$ is bounded
below by $\log2$, so that part of (NT-Perron-error) is
\[
 \ll x^cT^{-1}\sum_{n\ge1}d_D(n)n^{-c}
 =x^cT^{-1}\zeta(c)^D\ll xT^{-1}(\log x)^D.
\]
For $x/2<n<2x$, use $d_D(n)\ll_{D,\eta}x^\eta$ and
$|\log(x/n)|\gg |x-n|/x$.  Since $x$ is half-integral,
$\sum_{x/2<n<2x}|x-n|^{-1}\ll\log x$.  Also $(x/n)^c\ll1$,
uniformly for these $x,n$.  This part of the error is therefore
$O(x^{1+\eta}T^{-1}\log x)$.  Both parts are smaller than
$x^{1/2+\epsilon}$.

Move the vertical line to $\sigma_0=1/2+\delta$.  The integrand has no
pole in the intervening rectangle: $G$ is holomorphic there and
$\sigma_0>0$.  On either horizontal segment,
(NT-Perron-input) bounds the integral by
$O(x^cT^{\eta-1})=O(x^{-1+2\eta})$.
On the new vertical segment it gives
\[
 \ll x^{\sigma_0}\left(1+\int_1^Tt^{\eta-1}\,dt\right)
 \ll_\eta x^{1/2+\delta}T^\eta
 =O(x^{1/2+\epsilon/2}).
\]
The bounded part $|t|\le1$ is harmless because $\sigma_0>0$ and $G$
is bounded on its compact segment.  This proves the asserted partial-sum
bound with room in the exponent.

Apply this lemma first to $G(s)=1/L(s,\chi)$.
Its coefficients are $a_n=\mu(n)\chi(n)$, including the usual zero
values at integers not coprime to the conductor.  Under $A$, the
reciprocal is holomorphic on $\operatorname{Re}s>1/2$; the principal
L-function's pole produces a zero of the reciprocal, not a pole.
NT-L02 supplies (NT-Perron-input), so the twisted bound follows.
Character orthogonality, for $(a,q)=1$, is the finite identity
\[
 \sum_{\substack{n\le x\\n\equiv a\ ({\rm mod}\ q)}}\mu(n)
 =\frac1{\varphi(q)}\sum_{\chi\ ({\rm mod}\ q)}
       \overline{\chi(a)}\sum_{n\le x}\mu(n)\chi(n).
\]
Imprimitive functions differ by finitely many Euler factors whose zeros
lie on $\operatorname{Re}s=0$, and their reciprocals are bounded in any
closed half-strip to the right of $1/2$.  Thus each term satisfies the
same fixed-character argument.  A finite sum proves the fixed-$q$
progression bound, with no uniformity in $q$ claimed.

For a unitary Hecke character, expand
\[
 \frac1{L(s,\eta)}
 =\sum_{(\mathfrak a,\mathfrak f)=1}
         \mu_K(\mathfrak a)\eta(\mathfrak a)(N\mathfrak a)^{-s}
 =\sum_{n\ge1}a_nn^{-s}.
\]
For each $n$, $|a_n|$ is at most the number of ideals of norm $n$.
If $[K:\mathbb Q]=D$, comparison of the rational-prime Euler factors
of $\zeta_K$ with $(1-p^{-s})^{-D}$ gives this number at most $d_D(n)$.
Indeed there are at most $D$ prime ideals above $p$, each of residue
degree at least one, and replacing every local factor by a degree-one
factor can only increase its nonnegative coefficients.  Holomorphy and
(NT-Perron-input) follow conditionally on $M$ from NT-L03, including
pure norm twists.  The lemma proves NT-L10.  Averaging over the finite
character group of any fixed ray class group then projects onto one ray
class exactly as above; all sums are over ideals prime to its modulus.
This completes every claimed cancellation statement.

\paragraph{NT-Z10. The fixed infinite-order Hecke zero count.}
Let $\eta$ be any fixed unitary Hecke character of a number field $K$.
Its number of nontrivial zeros, counted with multiplicity, in
\[
 0\le\operatorname{Re}s\le1,\qquad |\operatorname{Im}s-T|\le1
\]
is $O_{K,\eta}(\log(|T|+3))$, unconditionally.  This assertion includes
infinite-order characters; its constant is not asserted uniform in
their infinity types.

Here is a full local argument.  On any fixed vertical strip the finite
Hecke L-function has a polynomial upper bound in $|t|$, away from its
possible pole.  In the present argument this follows directly from the
unconditional first part of Simoni\v c's Theorem~1, equation (2), on
$\sigma\ge1/2$, and from the functional equation and Stirling's formula
on the reflected strip.  The Selberg-class hypotheses for a non-pure-norm
character were verified in NT-L03 without using GRH; a pure norm twist
is a translate of the Dedekind function.  Consequently constants $C,A>0$
exist such that, for large $|T|$,
\[
 \max_{|s-(2+iT)|\le4}|L(s,\eta)|\le C(|T|+3)^A.
\tag{NT-J-upper}
\]
For a pure norm twist the only pole has a fixed imaginary coordinate, so
choose $|T|$ sufficiently large that this disk does not contain it.
In all other cases the function is entire.

The Euler product at the centre gives a lower bound that a general
Schwartz Mellin transform need not have:
\[
 |L(2+iT,\eta)|
 \ge\prod_{\mathfrak p\nmid\mathfrak f}(1+(N\mathfrak p)^{-2})^{-1}
 \ge\prod_{\mathfrak p}(1-(N\mathfrak p)^{-2})
 =\zeta_K(2)^{-1}>0.
\tag{NT-J-lower}
\]
These comparisons follow term by term from $|\eta(\mathfrak p)|=1$
and $1/(1+u)\ge1-u$ for $0\le u<1$, and persist on taking limits of
the absolutely convergent products.

Let $n_T(3)$ count zeros in the disk of radius $3$ about $2+iT$.
Jensen's formula on the radius-$4$ disk, or its limiting-radius version
if a zero lies on the outer circle, implies
\[
 n_T(3)\log(4/3)
 \le\log\max_{|s-(2+iT)|\le4}|L(s,\eta)|
                -\log|L(2+iT,\eta)|.
\]
The right side is $O_{K,\eta}(\log(|T|+3))$ by (NT-J-upper) and
(NT-J-lower).  Every point with $0\le\operatorname{Re}s\le1$ and
$|\operatorname{Im}s-T|\le1$ is at distance at most $\sqrt5<3$
from the centre, so its zeros are included in $n_T(3)$.  The remaining
bounded interval of $T$ meets only finitely many zeros, by meromorphic
analyticity and the absence of an identically zero function; increasing
the constant covers it.  This proves the claim.

The finite-order conductor-explicit counterpart is
Greni\'e--Molteni2019, Appendix~A, Theorem~A.1, author p.~26.
NT-Z10 instead proves the fixed-character form needed for the local
arithmetic divisor, and supplies the infinite-order version of root
MT-M04.  The argument uses the lower bound at its Jensen centre and
polynomial vertical growth.  It does not repair a statement quantified
over every function in an arbitrary Schwartz Mellin-transform space.

\subsection{An explicit uniform cyclotomic specialization}

\paragraph{NT-P11 and NT-P12.}
For $q\ge3$ and $(a,q)=1$, define
\[
 \psi(x;q,a)=\sum_{\substack{p^m\le x\\p^m\equiv a\ ({\rm mod}\ q)}}\log p,
 \qquad
 \pi(x;q,a)=\#\{p\le x:p\equiv a\ ({\rm mod}\ q)\}.
\]
Put $e_q=1$ if $q\equiv2\pmod4$ and $e_q=0$ otherwise.
Then $A$ implies
\begin{align*}
 \left|\psi(x;q,a)-\frac{x}{\varphi(q)}\right|
 &\le\sqrt{x}\left[
 \left(\frac{\log x}{2\pi}+2\right)\log q+
 \frac{\log^2x}{8\pi}+2\right]+e_q\log x,\qquad x\ge1,\\
 \left|\pi(x;q,a)-\frac{\operatorname{Li}_2(x)}{\varphi(q)}\right|
 &\le\sqrt{x}\left[
 \left(\frac1{2\pi}+\frac3{\log x}\right)\log q+
 \frac{\log x}{8\pi}+\frac1{4\pi}+\frac6{\log x}\right]+e_q,
 \qquad x\ge2.
\end{align*}

For the full proof, take $L=\mathbb Q(\zeta_q)$ and $K=\mathbb Q$.
The Galois group is $(\mathbb Z/q\mathbb Z)^\times$, of order
$n_L=\varphi(q)$, and its conjugacy classes are singletons.
The cyclotomic discriminant formula gives
\[
 \Delta_L=q^{\varphi(q)}
    \prod_{p\mid q}p^{-\varphi(q)/(p-1)},\qquad
 \log\Delta_L\le\varphi(q)\log q.
\]
At every prime $p\nmid q$, arithmetic Frobenius acts on $\zeta_q$ by
$\zeta_q\mapsto\zeta_q^p$.  Thus its $m$th power is the element
represented by $p^m$ modulo $q$.  The use of arithmetic or geometric
Frobenius merely replaces the labelled class by its inverse; the estimate
is identical after that relabelling.  The Dirichlet factorization of
$\zeta_L$ makes RH for it a consequence of $A$.
Apply NT-C01 and NT-C02 with $\delta_C=1/\varphi(q)$ and insert the
discriminant upper bound; all displayed coefficients multiplying
$\log\Delta_L$ are positive on the stated ranges.

If $q$ is odd or $4\mid q$, the primes unramified in $L$ are exactly
those not dividing $q$, so the resulting sums are the progression sums
just defined.  If $q\equiv2\pmod4$, write $q=2m$ with $m$ odd.
Then $L=\mathbb Q(\zeta_m)$ and the prime $2$ is unramified, whereas it
is excluded from a reduced class modulo $q$.  For odd primes the class
condition modulo $m$ is equivalent to that modulo $q$, since reduction
$(\mathbb Z/q\mathbb Z)^\times\to(\mathbb Z/m\mathbb Z)^\times$ is
an isomorphism.  The unramified Chebotarev sum can therefore differ only
by the powers of $2$ whose Frobenius powers lie in the designated class.
Their total weight is at most
$\lfloor\log x/\log2\rfloor\log2\le\log x$.
For the prime-counting sum the discrepancy is at most one.  Applying the
triangle inequality proves both formulas, including this degeneration.

The small moduli are transparent separately:
\begin{align*}
 \psi(x;1,0)&=\psi(x),& \pi(x;1,0)&=\pi(x),\\
 \psi(x;2,1)&=\psi(x)-\lfloor\log x/\log2\rfloor\log2,&
 \pi(x;2,1)&=\pi(x)-\mathbf1_{x\ge2}.
\end{align*}
For $1\le x<2$ the floor is zero.  Hence NT-P01 and NT-P03 give the
corresponding bounds in their stated ranges, with the explicit powers-of-two
correction retained for modulus two.  These are estimates on the open
cyclotomic arithmetic curve, and the additional term is exactly its
removed-prime contribution.

\subsection{A specified quadratic family and central nonvanishing}

\paragraph{NT-Z11. The exact one-level-density input.}
Let $\mathcal D$ consist of the nonzero odd squarefree integers, of both
signs.  For $d\in\mathcal D$, $8d$ is a fundamental discriminant and
$\chi_{8d}$ is its primitive quadratic character.  Put
$\ell_X=\log(X/(2\pi e))$.  Let $w$ be an even, nonnegative, nonzero
Schwartz function and write
\[
 W_w(X)=\sum_{d\in\mathcal D}w(d/X),\qquad
 \mathcal Z_{w,\phi}(X)=\frac1{W_w(X)}
  \sum_{d\in\mathcal D}w(d/X)
       \sum_{\rho=1/2+i\gamma}\phi(\gamma\ell_X/(2\pi)),
\]
where the inner sum counts every nontrivial zero with multiplicity.
Use the Fourier convention
$\widehat\phi(u)=\int_{\mathbb R}\phi(t)e^{-2\pi iut}\,dt$.
For a real even Schwartz function $\phi$ with
$\operatorname{supp}\widehat\phi\subset(-2,2)$, $A$ implies
\[
 \lim_{X\to\infty}\mathcal Z_{w,\phi}(X)
 =\widehat\phi(0)-\frac12\int_{-1}^1\widehat\phi(u)\,du.
\tag{NT-Z11}
\]
This is the leading term of Fiorilli--Parks--S\"odergren,
\texttt{arXiv:1601.06833}, Theorem~1.1, author p.~2, with definitions
(1.1)--(1.2) on the same page.  Its proof uses Theorem~3.5, author p.~12, and is concluded on p.~14.
The GRH input in the prime-sum step is Dirichlet GRH
(Lemma~2.5, author p.~7), which $A$ supplies.  No unrestricted Fourier
support, individual central nonvanishing, or nonabelian GRH is being
asserted here.

\paragraph{NT-Z12. A full nonvanishing deduction, including deweighting.}
Under $A$, at least $15/16$ of this specified family have a nonzero
central value in the following precise sense:
\[
 \liminf_{X\to\infty}
 \frac{\#\{d\in\mathcal D:|d|\le X,\ L(1/2,\chi_{8d})\ne0\}}
      {\#\{d\in\mathcal D:|d|\le X\}}\ge\frac{15}{16}.
\tag{NT-Z12}
\]
The same lower bound holds for the weighted proportions defined by any
of the fixed weights $w$ in NT-Z11.  This is a deduction from the cited
classical density theorem, not a claim of a new nonvanishing theorem.
Here are all approximation and counting details.

Let $m_d$ be the order of vanishing at $1/2$.  The completed primitive
quadratic function has root number $+1$.  One way to see the sign without
choosing a Gauss-sum square root is to divide the symmetric completed
Dedekind functional equation of $\mathbb Q(\sqrt{8d})$ by that of
$\mathbb Q$: the factorization
$\zeta_{\mathbb Q(\sqrt{8d})}=\zeta L(s,\chi_{8d})$ and the real or
complex gamma factors identify the quotient, up to a nonzero constant,
with the completed quadratic function.  The gamma duplication formula
handles the imaginary quadratic case.  Thus this entire function is
even about $s=1/2$.  Its gamma factor is finite and nonzero at the
centre, so $m_d$ is an even nonnegative integer.  In particular
\[
 2\mathbf1_{L(1/2,\chi_{8d})=0}\le m_d.
\tag{NT-central-parity}
\]
These functional equations are the standard Hecke equations; a precise
source for the completed convention is Ramakrishnan--Valenza,
Theorem~7-19, printed p.~279.

Choose real even smooth functions $g_j$ compactly supported in $(-1,1)$
such that $\int g_j=1$ and
$g_j\to g=\tfrac12\mathbf1_{[-1,1]}$ in $L^2(\mathbb R)$.
For example, smooth even cutoffs of $[-1+1/j,1-1/j]$, normalized by
their integrals, have these properties.  Define
\[
 f_j(t)=\int g_j(u)e^{2\pi iut}\,du,\qquad
 \phi_j(t)=|f_j(t)|^2.
\]
Then $\phi_j$ is real, even, nonnegative and Schwartz,
$\phi_j(0)=1$, and
$\operatorname{supp}\widehat\phi_j\subset(-2,2)$.
The support assertion follows because
$\widehat\phi_j=g_j*\widetilde g_j$, where
$\widetilde g_j(u)=\overline{g_j(-u)}$.
Plancherel gives $f_j\to f$ in $L^2$, with
\[
 f(t)=\frac{\sin(2\pi t)}{2\pi t}.
\]
Moreover, Cauchy--Schwarz yields
\[
 \|\phi_j-|f|^2\|_1
 \le\|f_j-f\|_2(\|f_j\|_2+\|f\|_2)\longrightarrow0.
\]
Consequently their Fourier transforms converge uniformly.  Since
\[
 \widehat{|f|^2}(u)=
 \begin{cases}\tfrac12-\tfrac14|u|,&|u|\le2,\\0,&|u|>2,
 \end{cases}
\]
we obtain
\[
 \lim_{j\to\infty}
 \left(\widehat\phi_j(0)-\frac12
                   \int_{-1}^1\widehat\phi_j(u)\,du\right)
 =\frac12-\frac12\cdot\frac34=\frac18.
\tag{NT-test-limit}
\]
Under $A$ every ordinate in NT-Z11 is real.  The nonnegativity of
$\phi_j$ and $\phi_j(0)=1$ therefore imply
\[
 \frac1{W_w(X)}\sum_{d\in\mathcal D}w(d/X)m_d
 \le\mathcal Z_{w,\phi_j}(X).
\]
First take $\limsup_{X\to\infty}$ for a fixed $j$, use NT-Z11, and
then let $j\to\infty$.  The result and (NT-central-parity) give
\[
 \limsup_{X\to\infty}
 \frac{\sum_{d\in\mathcal D}w(d/X)
               \mathbf1_{L(1/2,\chi_{8d})=0}}{W_w(X)}\le\frac1{16}.
\tag{NT-weighted-vanishing}
\]
This proves the weighted assertion without applying the density theorem
directly to a nonsmooth Fourier transform or to support touching $2$.

For completeness, the ordinary count follows by a separate elementary
argument.  If $A_+(Y)$ counts positive odd squarefree integers up to $Y$,
the squarefree indicator gives
\begin{align*}
 A_+(Y)&=\sum_{\substack{k\le\sqrt Y\\k\ {\rm odd}}}\mu(k)
       \#\{n\le Y/k^2:n\ {\rm odd}\}\\
 &=\frac Y2\sum_{k\ {\rm odd}}\frac{\mu(k)}{k^2}+O(\sqrt Y)
 =\frac4{\pi^2}Y+O(\sqrt Y).
\end{align*}
Here the omitted tail is $O(Y\sum_{k>\sqrt Y}k^{-2})=O(\sqrt Y)$,
and the Euler product is
$\sum_{k\ {\rm odd}}\mu(k)/k^2=8/\pi^2$.
Thus $N(X):=\#\{d\in\mathcal D:|d|\le X\}
 =8X/\pi^2+O(\sqrt X)$.
Partial summation of this counting formula gives, for every fixed
smooth compactly supported $w$,
\[
 W_w(X)=\frac4{\pi^2}X\int_{\mathbb R}w(t)\,dt+O_w(\sqrt X).
\]
Indeed apply Stieltjes integration to the positive count $A_+$ and
its reflected negative count; the integrated error is bounded by a
constant depending on $\int_0^\infty\sqrt t\,|w'(t)|dt$.

Fix $\delta>0$.  Take such an even smooth $w$ equal to $1$ on
$[-1,1]$, between $0$ and $1$ everywhere, and supported in
$[-1-\delta,1+\delta]$.  The number of vanishing members with
$|d|\le X$ is bounded above by its weighted counterpart.  Dividing by
$N(X)$ and using (NT-weighted-vanishing) gives a $\limsup$ of at most
$\frac1{16}\cdot\frac12\int w\le(1+\delta)/16$.
Letting $\delta\downarrow0$ proves (NT-Z12).

The characters in this statement are quadratic rank-one Artin objects.
Their varying conductors describe a family of double covers of open
arithmetic curves.  The symplectic expression in NT-Z11 matches the
finite-field hyperelliptic-family comparison discussed immediately after
Theorem~1.1 of the cited paper.  This comparison does not enlarge the
analytic support range or turn the $15/16$ bound into nonvanishing for
each member.  The ordinary proportion proved here concerns the
$\chi_{8d}$ family exactly; it has not been silently relabelled as all
fundamental discriminants.
